\backhead{OBJECTIVE TYPE QUESTIONS}{khred}
\begin{exlist}
\item The critical C:N ratio above which residues immobilise nitrogen is about\opt{5}{20}{80}{200}
\item The retardation factor of a non-sorbed solute is\opt{0}{1}{2}{$\rho_b/\theta$}
\item Denitrification is favoured by\opt{a dry soil}{a well-aerated soil}{a wet soil}{a frozen soil}
\item 50 mg NO$_3$\,L$^{-1}$ equals about\opt{3.5}{11.3}{25}{50} mg N\,L$^{-1}$
\item The half-life of a first-order process with $k=0.1$ d$^{-1}$ is\opt{0.69 d}{6.9 d}{10 d}{69 d}
\end{exlist}
\objkey{1.\ (b),\quad 2.\ (b),\quad 3.\ (c),\quad 4.\ (b),\quad 5.\ (b)}
\dykband{EXERCISES IN SWAFL}
{\small Open SWAFL (desktop or online) with the benchmark project. The chapter of the \href{https://github.com/halilunaziru73-creator/Coupled-Water-Fluxes-Pore-Scale-Dynamics-Precision-Irrigation-Mulching-Feedback-Loops/raw/main/Software/v1.1.0/SWAFL_Desktop/docs/manual_v1.1.0/SWAFL_Manual_v1.1.0.pdf}{SWAFL manual} that explains each task is given in brackets; record your values in a table and compare them with what is expected (\emph{Check}).}\par\vspace{3pt}
\begin{exlist}
\item {\color{khblue}[Manual, Chapter 22]} Open \emph{Nitrogen budget} and read the leaching of the wettest and the driest year.\par\noindent{\small\itshape Check: leaching is largest in the wettest year.}
\item {\color{khblue}[Manual, Chapter 28]} Design a Daisy experiment with three nitrogen rates and run it; compare the leaching.\par\noindent{\small\itshape Check: leaching rises with the rate.}
\item {\color{khblue}[Manual, Chapter 30]} Open the results of the chemical scenario and read the pesticide leached.\par\noindent{\small\itshape Check: most is degraded near the surface.}
\item {\color{khblue}[Manual, Chapter 25]} Click every box of the nitrogen part of the process map and list the processes.\par\noindent{\small\itshape Check: mineralisation, nitrification, denitrification, leaching, uptake.}
\end{exlist}
